% =========================================================================
\section{Seminarios prácticos}\label{sec:seminarios}
% =========================================================================
Esta sección resuelve, con todo detalle, los cinco problemas de la relación del tema~\cite{PROBLEMAS-T2}. Cada uno sigue la misma estructura que en
los temas anteriores: el \emph{enunciado}, la \emph{idea} que guía la resolución, la \emph{resolución} paso a paso y un breve comentario sobre la
\emph{interpretación} del resultado. Todos se reducen a calcular unos pocos momentos multipolares, de modo que la estrategia es siempre la misma:
\begin{enumerate}[label=(\arabic*)]
  \item elegir el origen (si el enunciado no lo fija, el centro de simetría de la distribución) y los ejes adaptados a la geometría;
  \item calcular $Q$, $\vp$ y $\Qjk$ con la \cref{tab:momentos}, aprovechando las simetrías para anular integrales antes de hacerlas
        (\cref{teo:revolucion}: con simetría de revolución basta $Q_{zz}$ y se puede partir directamente de \cref{eq:cuadrupolar-revolucion});
  \item decidir con la \cref{sec:error} cuántos términos hacen falta y construir el potencial; el campo se obtiene, si se pide, como $\vE=-\nab\phi$.
\end{enumerate}
La \cref{tab:resultados}, al final de la sección, recoge todos los resultados.

% =========================================================================
\subsection{Problema 2.1. Barra cargada vista desde lejos}\label{sem:barra}
% =========================================================================
\begin{ejercicio}[Potencial y campo de una barra con error menor del $1\%$]\label{ejer:barra}
Calcular, con un error menor del $1\%$, el campo eléctrico y el potencial creados por una distribución lineal de carga de densidad uniforme $\lambda$ sobre una línea
recta finita de longitud $\ell$, para puntos de observación situados a una distancia mayor que $10$ veces la longitud de la barra.
\end{ejercicio}

\begin{resolucion}
\textbf{Idea.} Tenemos una barra de longitud $\ell$ y densidad uniforme $\lambda$, y la observamos desde $r>10\ell$. No hace falta resolver la integral exacta: basta
desarrollar el potencial y quedarse con los términos que garantizan un error inferior al $1\%$. Tomamos el origen en el centro de la barra y el eje $z$ a lo largo de
ella (\figpanel{fig:sem-barra}{a}); así el problema tiene simetría de revolución alrededor de $z$ y simetría de inversión $z\to-z$.

\textbf{Paso 1: cuántos términos.} Toda la carga es del mismo signo, de modo que podemos usar la estimación del \cref{teo:cota-error}, en su versión con el tamaño
característico (\cref{eq:estimacion-media}): $\phimul{n}/\phimul{0}\lesssim(\langle r'\rangle/r)^{n}$. Para esta distribución, la longitud característica es del orden de la
propia barra, $\langle r'\rangle\sim\ell$, y como $r>10\ell$,
\[
  \frac{\langle r'\rangle}{r}\lesssim\frac{\ell}{r}<\frac{1}{10}\quad\Longrightarrow\quad\frac{\phimul{n}}{\phimul{0}}\lesssim\Bigl(\frac{1}{10}\Bigr)^{n}=10^{-n}.
\]
El término dipolar pesaría a lo sumo un $10\%$ y el cuadrupolar un $1\%$; para garantizar un error menor del $1\%$ basta calcular hasta el término cuadrupolar,
$\phi\simeq\phimul{0}+\phimul{1}+\phimul{2}$ (y veremos que $\phimul{1}=0$ por simetría).

\textbf{Paso 2: el desarrollo para la barra.} El elemento de carga es $\dif q=\lambda\,\dif z'$, en $\rp=z'\,\uz$ con $z'\in[-\ell/2,\ell/2]$. Si $z'>0$, el ángulo entre $\rp$ y
$\vc{r}$ es $\theta'=\theta$ y $r'=z'$; si $z'<0$, $\theta'=\pi-\theta$ y $r'=-z'$. Por la paridad de los polinomios de Legendre, $\Pl{n}(-x)=(-1)^{n}\Pl{n}(x)$, en ambos casos
$r'^{\,n}\Pl{n}(\cosp)=z'^{\,n}\Pl{n}(\cos\theta)$, y el \cref{teo:multipolar} da
\[
  \phi(\vc{r})=\frac{1}{4\pi\eo}\int_{-\ell/2}^{\ell/2}\frac{\lambda\,\dif z'}{\norm{\vc{r}-\rp}}
  =\frac{\lambda}{4\pi\eo}\sum_{n=0}^{\infty}\frac{\Pl{n}(\cos\theta)}{r^{n+1}}\int_{-\ell/2}^{\ell/2}z'^{\,n}\,\dif z' .
\]
\emph{Término monopolar} ($n=0$): $\int_{-\ell/2}^{\ell/2}\dif z'=\ell$ y $\Pl{0}=1$, luego
\[
  \phimul{0}(\vc{r})=\frac{\lambda\ell}{4\pi\eo r}=\frac{Q}{4\pi\eo r},\qquad Q=\lambda\ell .
\]
\emph{Término dipolar} ($n=1$): $\int_{-\ell/2}^{\ell/2}z'\,\dif z'=\bigl[\tfrac{z'^{2}}{2}\bigr]_{-\ell/2}^{\ell/2}=\tfrac{(\ell/2)^{2}}{2}-\tfrac{(\ell/2)^{2}}{2}=0$. Es consecuencia de haber
colocado el origen en el centro de la barra: el momento dipolar,
\[
  \vp=\int\rp\,\dif q=\lambda\int_{-\ell/2}^{\ell/2}z'\,\uz\,\dif z'=\lambda\,\uz\Bigl[\frac{z'^{2}}{2}\Bigr]_{-\ell/2}^{\ell/2}=\vc{0},
\]
se anula porque por cada elemento en $+z'$ hay otro en $-z'$ con la misma carga, $(+z')\lambda\,\dif z'+(-z')\lambda\,\dif z'=0$. Dicho con simetrías: la de revolución obliga a
que $\vp\parallel\uz$ (\cref{teo:revolucion}), y la de inversión $z\to-z$ a que $p_z=-p_z$, es decir, $\vp=\vc{0}$.

\emph{Término cuadrupolar} ($n=2$): $\int_{-\ell/2}^{\ell/2}z'^{2}\,\dif z'=\bigl[\tfrac{z'^{3}}{3}\bigr]_{-\ell/2}^{\ell/2}=\tfrac{\ell^{3}}{24}+\tfrac{\ell^{3}}{24}=\tfrac{\ell^{3}}{12}$ y
$\Pl{2}(\cos\theta)=\frac12(3\cos^{2}\theta-1)$, luego
\[
  \phimul{2}(\vc{r})=\frac{\lambda}{4\pi\eo r^{3}}\,\frac12\bigl(3\cos^{2}\theta-1\bigr)\frac{\ell^{3}}{12}=\frac{\lambda\ell^{3}}{96\pi\eo r^{3}}\bigl(3\cos^{2}\theta-1\bigr).
\]
\emph{Comprobación con el tensor cuadrupolar.} Como hay simetría de revolución, basta $Q_{zz}=\int(3z'^{2}-r'^{2})\,\dif q=\int2z'^{2}\,\lambda\,\dif z'=2\lambda\frac{\ell^{3}}{12}=\frac{\lambda\ell^{3}}{6}$,
y \cref{eq:cuadrupolar-revolucion} da $\phimul{2}=\frac{\lambda\ell^{3}/6}{16\pi\eo r^{3}}(3\cos^{2}\theta-1)$, el mismo resultado.

\textbf{Paso 3: el potencial.} Con un error menor del $1\%$,
\[
  \boxed{\;\phi(\vc{r})\simeq\frac{\lambda\ell}{4\pi\eo r}+\frac{\lambda\ell^{3}\bigl(3\cos^{2}\theta-1\bigr)}{96\pi\eo r^{3}}\qquad(r>10\ell)\;}
\]

\textbf{Paso 4: el campo.} El potencial solo depende de $r$ y $\theta$, de modo que, con el gradiente en esféricas,
$\vE=-\nab\phi=-\pd{\phi}{r}\,\ur-\frac1r\pd{\phi}{\theta}\,\uth$ ($E_\varphi=0$). Como hemos aproximado el potencial, el campo resultante es también aproximado.
\emph{Componente radial:}
\[
  E_r=-\pd{\phi}{r}=-\Bigl[-\frac{\lambda\ell}{4\pi\eo r^{2}}+\frac{\lambda\ell^{3}}{96\pi\eo}\bigl(3\cos^{2}\theta-1\bigr)(-3)r^{-4}\Bigr]
  =\frac{\lambda\ell}{4\pi\eo r^{2}}+\frac{\lambda\ell^{3}\bigl(3\cos^{2}\theta-1\bigr)}{32\pi\eo r^{4}} .
\]
\emph{Componente angular:} como $\pd{}{\theta}\bigl(3\cos^{2}\theta-1\bigr)=3\cdot2\cos\theta(-\sen\theta)=-6\sen\theta\cos\theta$,
\[
  \pd{\phi}{\theta}=\frac{\lambda\ell^{3}}{96\pi\eo r^{3}}\bigl(-6\sen\theta\cos\theta\bigr)=-\frac{\lambda\ell^{3}\sen\theta\cos\theta}{16\pi\eo r^{3}},\qquad
  E_\theta=-\frac1r\pd{\phi}{\theta}=\frac{\lambda\ell^{3}\sen\theta\cos\theta}{16\pi\eo r^{4}} .
\]
En definitiva,
\[
  \boxed{\;\vE(\vc{r})\simeq\Bigl[\frac{\lambda\ell}{4\pi\eo r^{2}}+\frac{\lambda\ell^{3}\bigl(3\cos^{2}\theta-1\bigr)}{32\pi\eo r^{4}}\Bigr]\ur+\frac{\lambda\ell^{3}\sen\theta\cos\theta}{16\pi\eo r^{4}}\,\uth\;}
\]
\end{resolucion}

% =========================================================================
% Figura del problema 2.1: (a) geometría de la barra (apuntes manuscritos,
% pág. 23) y (b) error relativo de la aproximación multipolar frente al
% potencial exacto. Datos: figuras/datos/error_barra.dat.
% =========================================================================
\begin{figure}[ht]
\centering
\begin{minipage}[c]{0.36\textwidth}\centering
\begin{tikzpicture}[x=1cm, y=1cm, >=Stealth]
  \draw[eje] (0,-1.9) -- (0,2.6) node[above, etiq] {$z$};
  \draw[eje] (0,0) -- (2.9,-0.55) node[right, etiq] {$y$};
  \draw[eje] (0,0) -- (-1.1,-0.95) node[below left=-2pt, etiq] {$x$};
  % barra
  \fill[magacero!85] (-0.06,-1.3) rectangle (0.06,1.3);
  \draw[fino] (0.15,1.3) -- (0.75,1.3) (0.15,-1.3) -- (0.75,-1.3);
  \draw[<->, fino] (0.6,-1.3) -- node[pos=0.22, right=1pt, etiq, font=\footnotesize] {$\ell$} (0.6,1.3);
  \node[etiq, font=\footnotesize, anchor=east] at (-0.12,1.45) {$\ell/2$};
  \node[etiq, font=\footnotesize, anchor=east] at (-0.12,-1.45) {$-\ell/2$};
  % elemento de carga
  \coordinate (Q) at (0,0.95);
  \fill[white, draw=magprof, line width=0.5pt] (-0.1,0.85) rectangle (0.1,1.05);
  \node[etiq, font=\footnotesize, anchor=east] at (-0.12,0.95) {$\dif q$};
  % punto campo
  \coordinate (O) at (0,0);
  \coordinate (Zp) at (0,1);
  \coordinate (P) at (2.4,2.45);
  \draw[vect] (O) -- (P) node[pos=0.55, below right=-1pt, etiq] {$\vc{r}$};
  \draw[->, line width=0.6pt, grispizarra, densely dashed] (Q) -- (P) node[pos=0.55, above=1pt, etiq] {$\vc{R}$};
  \pic[draw=magacero, line width=0.5pt, angle radius=0.42cm, "$\theta$", angle eccentricity=1.45, font=\small] {angle=P--O--Zp};
  \fill[magprof] (P) circle (1.5pt) node[right=2pt, etiq] {$P(r,\theta,\varphi)$};
\end{tikzpicture}\\[2pt]{\small\scshape (a)}
\end{minipage}\hfill
\begin{minipage}[c]{0.62\textwidth}\centering
\begin{tikzpicture}
\begin{loglogaxis}[marcoMMA, width=8.6cm, height=5.6cm, xmin=1, xmax=100, ymin=1e-9, ymax=40,
                   xlabel={$r/\ell$}, ylabel={error relativo (\%)},
                   every axis plot/.style={line width=1.05pt},
                   ytick={1e-8,1e-6,1e-4,1e-2,1}, xtick={1,10,100}, xticklabels={$1$,$10$,$100$},
                   legend pos=south west, legend cell align=left, legend style={font=\tiny}]
  \addplot[mmanaranja]                table[x=x, y=e0eje] {figuras/datos/error_barra.dat};
  \addplot[mmanaranja, densely dashed] table[x=x, y=e0ecu] {figuras/datos/error_barra.dat};
  \addplot[mmaazul]                   table[x=x, y=e2eje] {figuras/datos/error_barra.dat};
  \addplot[mmaazul, densely dashed]    table[x=x, y=e2ecu] {figuras/datos/error_barra.dat};
  \draw[rojoteo, line width=0.6pt] (axis cs:1,1) -- (axis cs:100,1);
  \draw[grispizarra, line width=0.6pt, densely dotted] (axis cs:10,1e-9) -- (axis cs:10,40);
  \node[font=\tiny, text=rojoteo, anchor=south east] at (axis cs:95,1) {$1\%$};
  \node[font=\tiny, text=grispizarra, anchor=south west, rotate=90, xshift=2pt] at (axis cs:10,3e-9) {$r=10\ell$};
  \legend{$\phimul{0}$; eje, $\phimul{0}$; ecuador, $\phimul{0}+\phimul{2}$; eje, $\phimul{0}+\phimul{2}$; ecuador}
\end{loglogaxis}
\end{tikzpicture}\\[2pt]{\small\scshape (b)}
\end{minipage}
\caption{Problema~2.1. (a)~Barra de longitud $\ell$ y densidad $\lambda$ sobre el eje $z$, centrada en el origen; el elemento de carga $\dif q=\lambda\,\dif z'$, en $z'$, y el punto campo $P$.
(b)~Error relativo, frente al potencial exacto, de la aproximación monopolar (naranja) y monopolar más cuadrupolar (azul), sobre el eje ($\theta=0$, continua) y en el plano
ecuatorial ($\theta=\pi/2$, a trazos). Para $r>10\ell$ (punteada) ambas aproximaciones quedan por debajo del $1\%$; la que incluye el término cuadrupolar mejora el resultado en
unos tres órdenes de magnitud.}
\label{fig:sem-barra}
\end{figure}


\begin{observacion}[Interpretación y un análisis más fino]\label{obs:sem-barra}
Lejos de la barra el término dominante es el de una carga puntual $Q=\lambda\ell$; la corrección cuadrupolar es positiva cerca del eje (la carga está «alargada» según $z$,
$Q_{zz}>0$) y negativa cerca del plano ecuatorial. La estimación $\langle r'\rangle\sim\ell$ es conservadora: la carga más alejada del centro está a $a=\ell/2$, de modo que
el \cref{teo:cota-error} garantiza $\abs{\phimul{2}/\phimul{0}}\le(1/20)^{2}=0{,}25\%$ y, de hecho, $\abs{\phimul{2}/\phimul{0}}=\frac{\ell^{2}}{24r^{2}}\abs{3\cos^{2}\theta-1}\le\frac{\ell^{2}}{12r^{2}}<0{,}084\%$.
Así pues, para $r>10\ell$ ya el término monopolar \emph{solo} cumple el requisito del $1\%$, y al añadir el cuadrupolar el error cae hasta el orden del término
siguiente no nulo, $\phimul{4}/\phimul{0}=\frac{\ell^{4}}{80r^{4}}\Pl{4}(\cos\theta)\sim10^{-6}$. La \figpanel{fig:sem-barra}{b} lo confirma comparando con el potencial exacto de la
barra, $\phi=\frac{\lambda}{4\pi\eo}\bigl[\operatorname{arcsenh}\frac{z+\ell/2}{s}-\operatorname{arcsenh}\frac{z-\ell/2}{s}\bigr]$ (con $s$ la distancia al eje).
\end{observacion}

% =========================================================================
\subsection{Problema 2.2. Momentos de una carga puntual}\label{sem:carga}
% =========================================================================
\begin{ejercicio}[Momentos de una carga puntual fuera del origen]\label{ejer:carga}
Sea una carga puntual $q$ situada en el punto $(a,b,c)$. Calcular el momento monopolar, dipolar y cuadrupolar para ese sistema. ¿Cuál de esas cantidades varía si se añade una
carga $-q$ en el origen de coordenadas?
\end{ejercicio}

\begin{resolucion}
\textbf{Idea.} Es un sistema discreto de una sola carga: los momentos se obtienen directamente de las expresiones discretas de la \cref{tab:momentos}, con
$\rp_1=a\,\ux+b\,\uy+c\,\uz$ y, por tanto, $r_1'^{\,2}=a^{2}+b^{2}+c^{2}$ (\figpanel{fig:sem-carga}{a}).

\textbf{1) Momentos del sistema inicial.} \emph{Monopolar} ($\ell=0$): la carga neta del sistema, $Q=\sum_iq_i=q$.

\emph{Dipolar} ($\ell=1$): el primer momento espacial de la distribución,
\[
  \vp=\sum_iq_i\,\rp_i=q\bigl(a\,\ux+b\,\uy+c\,\uz\bigr).
\]
\emph{Cuadrupolar} ($\ell=2$): $\Qjk=\sum_iq_i\bigl(3x'_{i,j}x'_{i,k}-\delt{j}{k}r_i'^{\,2}\bigr)$. En la diagonal principal ($j=k$, $\delt{j}{k}=1$),
\[
  Q_{xx}=q\bigl[3a^{2}-(a^{2}+b^{2}+c^{2})\bigr]=q\bigl(2a^{2}-b^{2}-c^{2}\bigr),
\]
\[
  Q_{yy}=q\bigl[3b^{2}-(a^{2}+b^{2}+c^{2})\bigr]=q\bigl(2b^{2}-a^{2}-c^{2}\bigr),\qquad
  Q_{zz}=q\bigl(2c^{2}-a^{2}-b^{2}\bigr).
\]
Fuera de la diagonal ($j\neq k$, $\delt{j}{k}=0$), y usando que el tensor es simétrico,
\[
  Q_{xy}=Q_{yx}=3qab,\qquad Q_{xz}=Q_{zx}=3qac,\qquad Q_{yz}=Q_{zy}=3qbc .
\]
En forma matricial,
\[
  \Qt=q\begin{pmatrix}2a^{2}-b^{2}-c^{2}&3ab&3ac\\3ab&2b^{2}-a^{2}-c^{2}&3bc\\3ac&3bc&2c^{2}-a^{2}-b^{2}\end{pmatrix},
\]
con $\Tr\Qt=q\,(2-1-1)(a^{2}+b^{2}+c^{2})=0$, como debe ser.

\textbf{2) Variaciones al añadir $-q$ en el origen} (\figpanel{fig:sem-carga}{b}). La nueva carga $q_2=-q$ tiene vector de posición nulo, $\rp_2=0\,\ux+0\,\uy+0\,\uz$.

\emph{Momento monopolar:} $Q_{\mathrm{nuevo}}=q_1+q_2=q+(-q)=0$. \emph{Es la única cantidad que varía} (y se anula).

\emph{Momentos dipolar y cuadrupolar:} sus contribuciones son proporcionales a las coordenadas de cada carga ($x'_j$) o a productos de ellas ($x'_jx'_k$, $r'^{2}$), y todas
se anulan para la carga situada en el origen:
\[
  \vp_{q_2}=(-q)\,\vc{0}=\vc{0},\qquad \bigl(\Qjk\bigr)_{q_2}=-q\bigl(3\cdot0\cdot0-\delt{j}{k}\cdot0\bigr)=0 .
\]
Por tanto, la carga añadida no perturba en absoluto los momentos de orden superior calculados antes.
\end{resolucion}

% =========================================================================
% Figura del problema 2.2: (a) carga q en (a,b,c); (b) tras añadir -q en el
% origen, el sistema es un dipolo de momento p = q r'_1. Perspectiva tdplot.
% =========================================================================
\begin{figure}[ht]
\centering
\newcommand{\ejesIII}{%
  \draw[eje] (0,0,0) -- (2.6,0,0) node[anchor=north east, etiq] {$x$};
  \draw[eje] (0,0,0) -- (0,2.7,0) node[anchor=west, etiq] {$y$};
  \draw[eje] (0,0,0) -- (0,0,2.3) node[anchor=south, etiq] {$z$};}
\newcommand{\proyecciones}{%
  \draw[punteado] (1.6,1.9,0) -- (1.6,1.9,1.5);
  \draw[punteado] (1.6,0,0) -- (1.6,1.9,0) -- (0,1.9,0);
  \node[etiq, font=\footnotesize, anchor=north] at (1.6,0,0) {$a$};
  \node[etiq, font=\footnotesize, anchor=north west] at (0,1.9,0) {$b$};
  \draw[punteado] (0,0,1.5) -- (1.6,1.9,1.5);
  \node[etiq, font=\footnotesize, anchor=east] at (0,0,1.5) {$c$};}
\tdplotsetmaincoords{70}{120}
\begin{minipage}[b]{0.47\textwidth}\centering
\begin{tikzpicture}[tdplot_main_coords, scale=1.25, >=Stealth]
  \ejesIII \proyecciones
  \draw[vect, rojoteo] (0,0,0) -- (1.6,1.9,1.5) node[pos=0.55, anchor=south east, etiq, text=rojoteo] {$\rp_1$};
  \node[cargapos] at (1.6,1.9,1.5) {$+$};
  \node[etiq, anchor=south west] at (1.6,1.9,1.62) {$q\ (a,b,c)$};
\end{tikzpicture}\\[2pt]{\small\scshape (a)}
\end{minipage}\hfill
\begin{minipage}[b]{0.47\textwidth}\centering
\begin{tikzpicture}[tdplot_main_coords, scale=1.25, >=Stealth]
  \ejesIII \proyecciones
  \draw[->, line width=1.5pt, magacero] (0,0,0) -- (1.6,1.9,1.5) node[pos=0.45, anchor=south east, xshift=-3pt, yshift=4pt, etiq, text=magacero] {$\vp=q\,\rp_1$};
  \node[cargapos] at (1.6,1.9,1.5) {$+$};
  \node[carganeg] at (0,0,0) {$-$};
  \node[etiq, anchor=south west] at (1.6,1.9,1.62) {$q$};
  \node[etiq, anchor=north west] at (0,0.1,-0.15) {$-q$};
\end{tikzpicture}\\[2pt]{\small\scshape (b)}
\end{minipage}
\caption{Problema~2.2. (a)~Carga puntual $q$ en $\rp_1=(a,b,c)$: $Q=q$, $\vp=q\,\rp_1$ y $\Qjk=q\,(3x'_{1,j}x'_{1,k}-\delt{j}{k}r_1'^{\,2})$. (b)~Al añadir $-q$ en el origen, el sistema se
convierte en un dipolo físico: solo cambia la carga neta ($Q=0$), mientras que $\vp$ y $\Qt$ no varían porque la nueva carga está en $\rp_2=\vc{0}$.}
\label{fig:sem-carga}
\end{figure}


\begin{observacion}[Interpretación]\label{obs:sem-carga}
Tras añadir la carga, el sistema es un dipolo físico de momento $\vp=q\,\rp_1$ y, al ser nula la carga neta, entra en juego el resultado de independencia del origen
(\cref{cor:centro-carga}): $\vp$ ya no cambia al trasladar los ejes. Antes de añadirla, en cambio, $\vp$ dependía del origen, y tomando el origen en la propia carga, que es su
centro de carga $\vc{r}_c=\vp/Q=(a,b,c)$, se anulaban $\vp$ y $\Qt$: una carga puntual vista desde su posición es un monopolo puro (\cref{ej:carga-origen}).
\end{observacion}

% =========================================================================
\subsection{Problema 2.3. Paralelepípedo uniformemente cargado}\label{sem:paralelepipedo}
% =========================================================================
\begin{ejercicio}[Potencial de un paralelepípedo con error menor del $1\%$]\label{ejer:paralelepipedo}
Sea una distribución de carga $Q$, uniforme, en un paralelepípedo de lados $2a$, $2b$, $2c$. Siendo todos los lados del mismo orden de magnitud $d$, calcular el potencial con un
error $<1\%$ cuando se observa ese recinto cargado a distancias mayores que $10d$.
\end{ejercicio}

\begin{resolucion}
\textbf{Idea.} El enunciado está incompleto: no indica respecto de qué origen, y como la carga neta no es nula el momento dipolar depende de esa elección
(\cref{prop:p-origen}). Lo natural es tomar el origen en el centro del paralelepípedo y los ejes paralelos a sus aristas (\cref{fig:sem-paralelepipedo}), de modo que ocupa la región
$[-a,a]\times[-b,b]\times[-c,c]$ con densidad $\rho=Q/\bigl[(2a)(2b)(2c)\bigr]=Q/(8abc)$.

\textbf{Paso 1: cuántos términos.} La carga es de un solo signo y su tamaño característico es $\langle r'\rangle\sim d$; con $r>10d$, $\phimul{n}/\phimul{0}\lesssim10^{-n}$, y para un
error menor del $1\%$ basta llegar al término cuadrupolar:
\[
  \phi\simeq\phimul{0}+\phimul{1}+\phimul{2}=\frac{Q}{4\pi\eo r}+\frac{\vp\cdot\ur}{4\pi\eo r^{2}}+\frac{1}{8\pi\eo r^{3}}\sum_{j,k}\Qjk\,\cd{j}\cd{k}.
\]

\textbf{Paso 2: momento dipolar.} Con $\rp=x'\,\ux+y'\,\uy+z'\,\uz$,
\[
  \vp=\rho\int_{-c}^{c}\!\!\int_{-b}^{b}\!\!\int_{-a}^{a}\bigl(x'\,\ux+y'\,\uy+z'\,\uz\bigr)\,\dif x'\,\dif y'\,\dif z'=\vc{0},
\]
porque cada componente es la integral de una función impar en un intervalo simétrico (por ejemplo, $\int_{-a}^{a}x'\,\dif x'=0$).

\textbf{Paso 3: tensor cuadrupolar.} En la diagonal, $3x'^{2}-r'^{2}=2x'^{2}-y'^{2}-z'^{2}$, y como las integrales se separan,
\begin{align*}
  Q_{xx}&=\rho\int_{-c}^{c}\!\!\int_{-b}^{b}\!\!\int_{-a}^{a}\bigl(2x'^{2}-y'^{2}-z'^{2}\bigr)\dif x'\,\dif y'\,\dif z'\\
         &=\rho\Bigl[2\cdot2b\cdot2c\Bigl[\frac{x'^{3}}{3}\Bigr]_{-a}^{a}-2a\cdot2c\Bigl[\frac{y'^{3}}{3}\Bigr]_{-b}^{b}-2a\cdot2b\Bigl[\frac{z'^{3}}{3}\Bigr]_{-c}^{c}\Bigr]\\
        &=\rho\Bigl[2\cdot4bc\,\frac{2a^{3}}{3}-4ac\,\frac{2b^{3}}{3}-4ab\,\frac{2c^{3}}{3}\Bigr]\\
        &=\underbrace{8\rho\,abc}_{Q}\,\frac{2a^{2}-b^{2}-c^{2}}{3}=\frac{Q}{3}\bigl(2a^{2}-b^{2}-c^{2}\bigr),
\end{align*}
y, de forma análoga (permutando los papeles de $a$, $b$, $c$),
\[
  Q_{yy}=\frac{Q}{3}\bigl(2b^{2}-a^{2}-c^{2}\bigr),\qquad Q_{zz}=\frac{Q}{3}\bigl(2c^{2}-a^{2}-b^{2}\bigr).
\]
Fuera de la diagonal, $Q_{xy}=3\rho\iiint x'y'\,\dif x'\,\dif y'\,\dif z'=0$ (integrando impar en $x'$), e igual $Q_{xz}=Q_{yz}=0$: los ejes de las aristas son los ejes principales.
Comprobación: $\Tr\Qt=\frac{Q}{3}\bigl[(2-1-1)a^{2}+(2-1-1)b^{2}+(2-1-1)c^{2}\bigr]=0$.

\textbf{Paso 4: el potencial.} Sustituyendo,
\[
  \phi(\vc{r})\simeq\frac{Q}{4\pi\eo r}+\frac{Q}{24\pi\eo r^{3}}\Bigl[\bigl(2a^{2}-b^{2}-c^{2}\bigr)\cd{x}^{2}+\bigl(2b^{2}-a^{2}-c^{2}\bigr)\cd{y}^{2}+\bigl(2c^{2}-a^{2}-b^{2}\bigr)\cd{z}^{2}\Bigr],
\]
con $\cd{x}^{2}=\sen^{2}\theta\cos^{2}\varphi=x^{2}/r^{2}$, $\cd{y}^{2}=\sen^{2}\theta\sen^{2}\varphi=y^{2}/r^{2}$ y $\cd{z}^{2}=\cos^{2}\theta=z^{2}/r^{2}$. En cartesianas se aprecia mejor la simetría:
\[
  \boxed{\;\phi(\vc{r})\simeq\frac{Q}{4\pi\eo r}+\frac{Q}{24\pi\eo r^{5}}\Bigl[x^{2}\bigl(2a^{2}-b^{2}-c^{2}\bigr)+y^{2}\bigl(2b^{2}-a^{2}-c^{2}\bigr)+z^{2}\bigl(2c^{2}-a^{2}-b^{2}\bigr)\Bigr]\;}
\]
\end{resolucion}

% =========================================================================
% Figura del problema 2.3 (apuntes manuscritos, pág. 19): paralelepípedo
% [-a,a]x[-b,b]x[-c,c] uniformemente cargado, con el origen en su centro.
% =========================================================================
\begin{figure}[ht]
\centering
\tdplotsetmaincoords{68}{122}
\begin{tikzpicture}[tdplot_main_coords, scale=1.05, >=Stealth]
  \def\A{1.5}\def\B{2.1}\def\C{1.0}
  % caras ocultas (aristas a trazos)
  \draw[magacero!70, densely dashed, line width=0.5pt] (-\A,-\B,-\C) -- (\A,-\B,-\C) (-\A,-\B,-\C) -- (-\A,\B,-\C) (-\A,-\B,-\C) -- (-\A,-\B,\C);
  % caras visibles
  \fill[rosaclaro, opacity=0.55] (\A,-\B,-\C) -- (\A,\B,-\C) -- (\A,\B,\C) -- (\A,-\B,\C) -- cycle;
  \fill[rosaclaro!70, opacity=0.55] (-\A,\B,-\C) -- (\A,\B,-\C) -- (\A,\B,\C) -- (-\A,\B,\C) -- cycle;
  \fill[rosaclaro!45, opacity=0.6] (-\A,-\B,\C) -- (\A,-\B,\C) -- (\A,\B,\C) -- (-\A,\B,\C) -- cycle;
  \draw[magacero, line width=0.7pt] (\A,-\B,-\C) -- (\A,\B,-\C) -- (-\A,\B,-\C) -- (-\A,\B,\C) -- (-\A,-\B,\C) -- (\A,-\B,\C) -- cycle;
  \draw[magacero, line width=0.7pt] (\A,\B,-\C) -- (\A,\B,\C) -- (-\A,\B,\C) (\A,\B,\C) -- (\A,-\B,\C);
  % ejes por el centro
  \draw[eje] (0,0,0) -- (2.9,0,0) node[anchor=north east, etiq] {$x$};
  \draw[eje] (0,0,0) -- (0,3.3,0) node[anchor=west, etiq] {$y$};
  \draw[eje] (0,0,0) -- (0,0,2.2) node[anchor=south, etiq] {$z$};
  % elemento de volumen
  \coordinate (Q) at (1.0,0.6,0.3);
  \draw[vect] (0,0,0) -- (Q) node[pos=0.55, anchor=north west, etiq] {$\rp$};
  \fill[white, draw=magprof, line width=0.5pt] (Q) circle (2pt);
  \node[etiq, font=\footnotesize, anchor=south west] at ($(Q)+(0,0.08,0.05)$) {$\dVp$};
  % cotas
  \draw[<->, fino] (\A+0.25,-\B,-\C) -- node[anchor=north, etiq, font=\footnotesize] {$2b$} (\A+0.25,\B,-\C);
  \draw[<->, fino] (\A,\B+0.25,-\C) -- node[anchor=north west, etiq, font=\footnotesize] {$2a$} (-\A,\B+0.25,-\C);
  \draw[<->, fino] (\A,\B+0.25,-\C) -- node[anchor=west, etiq, font=\footnotesize] {$2c$} (\A,\B+0.25,\C);
  \node[etiq, text=magprof, anchor=south east] at (-\A,-\B,\C) {$\rho=\dfrac{Q}{8abc}$};
\end{tikzpicture}
\caption{Problema~2.3. Paralelepípedo de lados $2a$, $2b$, $2c$ con carga $Q$ uniforme, centrado en el origen y con las aristas paralelas a los ejes. La simetría de inversión anula
el momento dipolar y todas las componentes no diagonales del tensor cuadrupolar.}
\label{fig:sem-paralelepipedo}
\end{figure}


\begin{observacion}[Interpretación: el cubo y la paridad]\label{obs:sem-paralelepipedo}
\begin{enumerate}[label=(\alph*)]
  \item La corrección cuadrupolar es positiva en la dirección de la arista más larga (si $a>b,c$, $Q_{xx}>0$): la caja se comporta como una carga «alargada» según ese eje.
  \item Para un \emph{cubo}, $a=b=c$, todas las componentes se anulan, $\Qt=0$: la simetría cúbica es suficiente para que el potencial no tenga corrección cuadrupolar, aunque la
        distribución no sea esférica.
  \item La distribución es simétrica frente a la inversión $\rp\to-\rp$, de modo que \emph{todos} los momentos de orden impar se anulan (sus integrandos, $r'^{\,\ell}\Pl{\ell}(\cosp)$
        con $\ell$ impar, cambian de signo). El primer término despreciado no es, pues, el octupolar sino el hexadecapolar ($\ell=4$), y el error real es del orden de $(d/r)^{4}\sim10^{-4}$,
        muy por debajo del $1\%$ exigido.
\end{enumerate}
\end{observacion}

% =========================================================================
\subsection{Problema 2.4. Esfera con densidad superficial \texorpdfstring{$\sigma_0\cos\theta$}{σ0 cos θ}}\label{sem:esfera-sigma}
% =========================================================================
\begin{ejercicio}[Momento dipolar de una esfera polarizada]\label{ejer:esfera-sigma}
Calcular el momento dipolar de una distribución superficial de carga sobre una esfera de radio $a$ cuando la densidad varía con el ángulo $\theta$ de la forma
$\sigma(\theta)=\sigma_0\cos\theta$.
\end{ejercicio}

\begin{resolucion}
\textbf{Idea.} La carga es positiva en el hemisferio norte y negativa en el sur (\figpanel{fig:sem-esfera-sigma}{a}): esperamos un momento dipolar dirigido según $+\uz$. Antes de
calcularlo comprobamos que la carga neta es nula, porque entonces el resultado no depende del origen y podemos tomar como referencia el centro de la esfera:
con $\dSp=a^{2}\sen\theta'\,\dif\theta'\,\dif\varphi'$,
\[
  Q=\oiint_{S'}\sigma\,\dSp=\sigma_0a^{2}\int_{0}^{2\pi}\dif\varphi'\int_{0}^{\pi}\cos\theta'\sen\theta'\,\dif\theta'=\sigma_0a^{2}\cdot2\pi\cdot\Bigl[\frac{\sen^{2}\theta'}{2}\Bigr]_{0}^{\pi}=0 .
\]

\textbf{Paso 1: el momento dipolar de una distribución superficial.} El momento dipolar de una distribución volúmica es $\vp=\int_{V'}\rp\rho(\rp)\,\dVp$. Al pasar a una distribución
estrictamente bidimensional sobre la superficie $S'$, sustituimos la densidad volúmica por la superficial y el diferencial de volumen por el de área:
\[
  \vp=\oiint_{S'}\rp\,\sigma(\theta')\,\dSp .
\]
\emph{Justificación mediante la delta de Dirac.} Una distribución superficial no es más que el caso límite de una volúmica en la que toda la carga está confinada en una capa de espesor
nulo: en esféricas, $\rho(r',\theta',\varphi')=\sigma(\theta',\varphi')\,\delta(r'-a)$, dimensionalmente correcta porque $[\mathrm{C/m^{3}}]=[\mathrm{C/m^{2}}][\mathrm{m^{-1}}]$.
Sustituyendo, con $\rp=r'\,\ur$,
\[
  \vp=\int_{0}^{2\pi}\!\!\int_{0}^{\pi}\!\!\int_{0}^{\infty}\bigl(r'\,\ur\bigr)\,\sigma(\theta',\varphi')\,\delta(r'-a)\,r'^{2}\sen\theta'\,\dif r'\,\dif\theta'\,\dif\varphi' .
\]
Reagrupando los factores radiales, la integral en $r'$ es $\int_{0}^{\infty}r'^{3}\delta(r'-a)\,\dif r'=a^{3}$, de modo que
$\vp=\int_{0}^{2\pi}\!\int_{0}^{\pi}(a\,\ur)\,\sigma\,\bigl(a^{2}\sen\theta'\,\dif\theta'\,\dif\varphi'\bigr)=\oiint_{S'}\rp\big|_{r'=a}\,\sigma\,\dSp$, como queríamos.

\textbf{Paso 2: parametrización.} Esta integral es \emph{vectorial} (no es un flujo $\iint\vc{F}\cdot\un\,\dsc$): cada componente cartesiana es una integral escalar. Usamos
\[
  \rp=a\,\ur=a\bigl(\sen\theta'\cos\varphi'\,\ux+\sen\theta'\sen\varphi'\,\uy+\cos\theta'\,\uz\bigr),\qquad \dSp=a^{2}\sen\theta'\,\dif\theta'\,\dif\varphi' .
\]

\textbf{Paso 3: componentes transversales.} Por el teorema de Fubini (producto de integrales simples iteradas),
\begin{align*}
  p_x&=\int_{0}^{2\pi}\!\!\int_{0}^{\pi}\bigl(a\sen\theta'\cos\varphi'\bigr)\bigl(\sigma_0\cos\theta'\bigr)\bigl(a^{2}\sen\theta'\bigr)\dif\theta'\,\dif\varphi'\\
  &=a^{3}\sigma_0\int_{0}^{\pi}\sen^{2}\theta'\cos\theta'\,\dif\theta'\underbrace{\int_{0}^{2\pi}\cos\varphi'\,\dif\varphi'}_{[\sen\varphi']_{0}^{2\pi}=0}=0,
\end{align*}
\[
  p_y=a^{3}\sigma_0\int_{0}^{\pi}\sen^{2}\theta'\cos\theta'\,\dif\theta'\underbrace{\int_{0}^{2\pi}\sen\varphi'\,\dif\varphi'}_{[-\cos\varphi']_{0}^{2\pi}=0}=0 .
\]
Las componentes transversales al eje de revolución se anulan idénticamente porque $\sigma=\sigma_0\cos\theta'$ tiene simetría azimutal perfecta (no depende de $\varphi'$),
como asegura el \cref{teo:revolucion}.

\textbf{Paso 4: componente axial.} Separando las variables angulares,
\[
  \vp=p_z\,\uz=a^{3}\sigma_0\,\uz\int_{0}^{2\pi}\dif\varphi'\int_{0}^{\pi}\cos^{2}\theta'\sen\theta'\,\dif\theta' .
\]
Con el cambio $u=\cos\theta'$, $\dif u=-\sen\theta'\,\dif\theta'$, que lleva $\theta'=0\mapsto u=1$ y $\theta'=\pi\mapsto u=-1$,
\[
  \int_{0}^{\pi}\cos^{2}\theta'\sen\theta'\,\dif\theta'=-\int_{1}^{-1}u^{2}\,\dif u=\int_{-1}^{1}u^{2}\,\dif u=\Bigl[\frac{u^{3}}{3}\Bigr]_{-1}^{1}=\frac13-\Bigl(-\frac13\Bigr)=\frac23 .
\]
Por tanto,
\[
  \boxed{\;\vp=a^{3}\sigma_0\,(2\pi)\,\frac23\,\uz=\frac43\pi a^{3}\sigma_0\,\uz\;}
\]
\end{resolucion}

% =========================================================================
% Figura del problema 2.4 (apuntes manuscritos, pág. 19): (a) esfera con
% sigma = sigma_0 cos(theta) (signos con densidad proporcional a |cos t|) y
% su momento dipolar; (b) líneas de campo exactas (analíticas): dipolares,
% r = C sen^2(t), fuera, y uniformes (verticales, hacia -z) dentro, con las
% equipotenciales a trazos (r^2 = a^2 cos t / cos t1 fuera, z = cte dentro).
% =========================================================================
\begin{figure}[ht]
\centering
\begin{minipage}[b]{0.45\textwidth}\centering
\begin{tikzpicture}[x=1cm, y=1cm, >=Stealth]
  \def\ra{1.55}
  \draw[eje] (0,-2.1) -- (0,2.4) node[above, etiq] {$z$};
  \draw[eje] (-2.1,0) -- (2.4,0) node[right, etiq] {$x$};
  \shade[inner color=white, outer color=rosaclaro] (0,0) circle (\ra);
  \draw[magacero, line width=0.8pt] (0,0) circle (\ra);
  % signos: más densos cerca de los polos (|sigma| ~ |cos t|)
  \foreach \t in {8,22,36,50,64}{%
    \foreach \s in {1,-1}{%
      \pgfmathsetmacro{\tam}{0.35+0.65*cos(\t)}%
      \node[font=\scriptsize\bfseries, text=rojoteo, scale=\tam] at ({\s*(\ra+0.22)*sin(\t)},{(\ra+0.22)*cos(\t)}) {$+$};
      \node[font=\scriptsize\bfseries, text=pizarra, scale=\tam] at ({\s*(\ra+0.22)*sin(\t)},{-(\ra+0.22)*cos(\t)}) {$-$};}}
  % momento dipolar
  \draw[->, line width=1.6pt, magprof] (0,-0.55) -- (0,0.75) node[left=1pt, etiq, text=magprof] {$\vp$};
  % punto fuente y ángulo
  \coordinate (O) at (0,0);
  \coordinate (Zp) at (0,1);
  \coordinate (S) at ({\ra*sin(52)},{\ra*cos(52)});
  \draw[vect, rojoteo] (O) -- (S);
  \node[etiq, text=rojoteo, font=\footnotesize, anchor=west] at (1.74,0.55) {$\rp=a\,\ur$};
  \draw[rojoteo!60, line width=0.3pt] (1.76,0.55) -- (0.72,0.52);
  \pic[draw=magacero, line width=0.5pt, angle radius=0.62cm, "$\theta'$", angle eccentricity=1.35, font=\scriptsize] {angle=S--O--Zp};
  \fill[magprof] (S) circle (1.3pt) node[above right=0pt, etiq, font=\footnotesize] {$\dSp$};
\end{tikzpicture}\\[2pt]{\small\scshape (a)}
\end{minipage}\hfill
\begin{minipage}[b]{0.5\textwidth}\centering
\begin{tikzpicture}[x=0.95cm, y=0.95cm, >=Stealth]
  \def\ra{1.2}
  \begin{scope}
    \clip (-3,-2.6) rectangle (3,2.6);
    % equipotenciales (a trazos)
    \foreach \tu in {40,70}{%
      \foreach \s in {1,-1}{%
        \draw[lequi] ({-\ra*sin(\tu)},{\s*\ra*cos(\tu)}) -- ({\ra*sin(\tu)},{\s*\ra*cos(\tu)});
        \draw[lequi, domain=0:89, samples=60, variable=\t]
          plot ({\ra*sqrt(cos(\t)/cos(\tu))*sin(\t)},{\s*\ra*sqrt(cos(\t)/cos(\tu))*cos(\t)});
        \draw[lequi, domain=0:89, samples=60, variable=\t]
          plot ({-\ra*sqrt(cos(\t)/cos(\tu))*sin(\t)},{\s*\ra*sqrt(cos(\t)/cos(\tu))*cos(\t)});}}
    \draw[lequi] (-3,0) -- (3,0);
    % líneas de campo exteriores: r = a sen^2(t)/sen^2(t0), de t0 a 180-t0
    \foreach \tc in {22,38,56}{%
      \pgfmathsetmacro{\C}{\ra/(sin(\tc)^2)}%
      \foreach \s in {1,-1}{%
        \draw[lcampo, domain=\tc:180-\tc, samples=90, variable=\t]
          plot ({\s*\C*sin(\t)^3},{\C*sin(\t)^2*cos(\t)});}}
    % líneas de campo interiores: verticales, hacia -z
    \foreach \tc in {22,38,56}{%
      \foreach \s in {1,-1}{%
        \draw[lcampo] ({\s*\ra*sin(\tc)},{\ra*cos(\tc)}) -- ({\s*\ra*sin(\tc)},{-\ra*cos(\tc)});}}
    \draw[lcampo] (0,\ra) -- (0,-\ra);
    \draw[lcampoinv] (0,2.6) -- (0,\ra);
    \draw[lcampo] (0,-\ra) -- (0,-2.6);
  \end{scope}
  \draw[magprof, line width=0.8pt] (0,0) circle (\ra);
  \draw[grispizarra!60, line width=0.5pt] (-3,-2.6) rectangle (3,2.6);
  \node[etiq, font=\scriptsize, fill=white, inner sep=0.8pt] at (0.25,-0.32) {$\vE_{\mathrm{int}}$};
\end{tikzpicture}\\[2pt]{\small\scshape (b)}
\end{minipage}
\caption{Problema~2.4. (a)~Esfera de radio $a$ con densidad superficial $\sigma=\sigma_0\cos\theta$: positiva en el hemisferio norte y negativa en el sur (el tamaño de cada signo
crece con $\abs{\sigma}$), y su momento dipolar $\vp=\frac43\pi a^{3}\sigma_0\,\uz$. (b)~Líneas de campo exactas (magenta) y equipotenciales (a trazos) en un plano meridiano: fuera de
la esfera el campo es el de un dipolo puntual en el centro, $r=C\sen^{2}\theta$; dentro, uniforme y dirigido según $-\uz$. Las líneas nacen en el hemisferio positivo y mueren en
el negativo, tanto por fuera como por dentro.}
\label{fig:sem-esfera-sigma}
\end{figure}


\begin{observacion}[Interpretación: un dipolo exacto]\label{obs:sem-esfera-sigma}
El momento es el de una esfera de volumen $\frac43\pi a^{3}$ con polarización uniforme $\vP=\sigma_0\,\uz$: en efecto, por el \cref{teo:cargas-pol}, una esfera uniformemente polarizada
tiene $\rhoP=-\nabp\cdot\vP=0$ y $\sigP=\vP\cdot\ur=\sigma_0\cos\theta$, que es justamente nuestra distribución. Más aún, como $\sigma\propto\Pl{1}(\cos\theta')$, la ortogonalidad de los
polinomios de Legendre hace que todos los momentos con $\ell\neq1$ se anulen, y el potencial es \emph{exactamente} dipolar fuera de la esfera y el de un campo uniforme
dentro~\cite{GRIFFITHS}:
\[
  \phi(\vc{r})=\begin{cases}\dfrac{\sigma_0}{3\eo}\,r\cos\theta=\dfrac{\sigma_0}{3\eo}\,z,& r\le a,\\[10pt]
  \dfrac{\sigma_0a^{3}}{3\eo}\,\dfrac{\cos\theta}{r^{2}}=\dfrac{\vp\cdot\ur}{4\pi\eo r^{2}},& r\ge a,\end{cases}
  \qquad \vE_{\mathrm{int}}=-\frac{\sigma_0}{3\eo}\,\uz .
\]
Ambas expresiones coinciden en $r=a$, como exige la continuidad del potencial~\cite{TEMA1}, y la \figpanel{fig:sem-esfera-sigma}{b} muestra las líneas de campo
correspondientes: dipolares fuera y paralelas dentro.
\end{observacion}

% =========================================================================
\subsection{Problema 2.5. Distribución con simetría esférica}\label{sem:esfera-rho}
% =========================================================================
\begin{ejercicio}[Momentos de una distribución esféricamente simétrica]\label{ejer:esfera-rho}
Demostrar que una distribución de carga esférica y esféricamente simétrica da lugar a momentos dipolar y cuadrupolar nulos cuando se elige un sistema de coordenadas centrado en la
esfera.
\end{ejercicio}

\begin{resolucion}
\textbf{Idea.} En una distribución esféricamente simétrica centrada en el origen la densidad de carga solo depende de la coordenada radial, $\rho=\rho(r')$. Sea la carga confinada en una
esfera de radio $R$ (\cref{fig:sem-esfera-rho}). Expresamos el vector de posición del punto fuente y el elemento de volumen en esféricas,
\[
  \rp=r'\bigl(\sen\theta'\cos\varphi'\,\ux+\sen\theta'\sen\varphi'\,\uy+\cos\theta'\,\uz\bigr),\qquad \dVp=r'^{2}\sen\theta'\,\dif r'\,\dif\theta'\,\dif\varphi',
\]
y aprovechamos que todas las integrales se separan en un factor radial y dos angulares.

\textbf{Momento dipolar.} Separando $\vp=\int_{V'}\rp\rho(r')\,\dVp$ en sus componentes cartesianas,
\[
  p_x=\Bigl[\int_{0}^{R}r'^{3}\rho(r')\,\dif r'\Bigr]\Bigl[\int_{0}^{\pi}\sen^{2}\theta'\,\dif\theta'\Bigr]\Bigl[\int_{0}^{2\pi}\cos\varphi'\,\dif\varphi'\Bigr]=0,
\]
porque la integral azimutal es $[\sen\varphi']_{0}^{2\pi}=0$ (es nula al ser una función sinusoidal integrada en un ciclo completo). Por la misma razón,
$p_y\propto\int_{0}^{2\pi}\sen\varphi'\,\dif\varphi'=0$. En la componente $z$ es la integral \emph{polar} la que se anula:
\[
  p_z=\Bigl[\int_{0}^{R}r'^{3}\rho\,\dif r'\Bigr]\Bigl[\int_{0}^{\pi}\cos\theta'\sen\theta'\,\dif\theta'\Bigr]\Bigl[\int_{0}^{2\pi}\dif\varphi'\Bigr],
\]
\[
  \int_{0}^{\pi}\cos\theta'\sen\theta'\,\dif\theta'=\Bigl[\frac{\sen^{2}\theta'}{2}\Bigr]_{0}^{\pi}=\frac12\bigl(\sen^{2}\pi-\sen^{2}0\bigr)=0,
\]
donde hemos usado $\frac{\dif}{\dif\theta'}\bigl(\sen^{2}\theta'\bigr)=2\sen\theta'\cos\theta'$ y el teorema fundamental del cálculo. Concluimos que $\vp=\vc{0}$.

\textbf{Momento cuadrupolar.} Una distribución esféricamente simétrica tiene, en particular, simetría de revolución alrededor de $z$, así que por el \cref{teo:revolucion} (dos
ligaduras: traza nula y $Q_{xx}=Q_{yy}$)
\[
  \Qjk=0\ (j\neq k),\qquad Q_{xx}=Q_{yy}=-\tfrac12Q_{zz}.
\]
(Se comprueba también directamente: por ejemplo, $Q_{xy}=3\int r'^{2}\sen^{2}\theta'\cos\varphi'\sen\varphi'\,\rho\,\dVp\propto\int_0^{2\pi}\cos\varphi'\sen\varphi'\,\dif\varphi'=0$.) Basta,
pues, demostrar que $Q_{zz}=0$. Con $z'=r'\cos\theta'$ y $\delt{z}{z}=1$,
\[
  Q_{zz}=\int_{V'}\bigl(3z'^{2}-r'^{2}\bigr)\rho\,\dVp=\Bigl[\int_{0}^{R}r'^{4}\rho(r')\,\dif r'\Bigr]\Bigl[\int_{0}^{\pi}\bigl(3\cos^{2}\theta'-1\bigr)\sen\theta'\,\dif\theta'\Bigr]\Bigl[\int_{0}^{2\pi}\dif\varphi'\Bigr].
\]
El factor radial es finito y el azimutal vale $2\pi$; todo se decide en el factor polar. Con $u=\cos\theta'$,
\[
  \int_{0}^{\pi}\bigl(3\cos^{2}\theta'-1\bigr)\sen\theta'\,\dif\theta'=\int_{-1}^{1}\bigl(3u^{2}-1\bigr)\dif u=\bigl[u^{3}-u\bigr]_{-1}^{1}=0-0=0 .
\]
Concluimos que $Q_{zz}=0$ y, con ello, $Q_{xx}=Q_{yy}=-\frac12Q_{zz}=0$ y $\Qjk=0$ para todo $j,k$.
\end{resolucion}

% =========================================================================
% Figura del problema 2.5 (apuntes manuscritos, pág. 27): distribución con
% simetría esférica rho(r) confinada en una esfera de radio R (sombreado
% radial) y geometría del cálculo de los momentos.
% =========================================================================
\begin{figure}[ht]
\centering
\begin{tikzpicture}[x=1cm, y=1cm, >=Stealth]
  \def\Rr{1.7}
  \shade[inner color=magacero!75, outer color=rosaclaro!60] (0,0) circle (\Rr);
  \draw[magacero, line width=0.8pt] (0,0) circle (\Rr);
  \draw[magacero!60, line width=0.4pt, densely dashed] (0,0) ellipse ({\Rr} and {0.42*\Rr});
  \draw[eje] (0,0) -- (0,2.5) node[above, etiq] {$z$};
  \draw[eje] (0,0) -- (2.6,-0.35) node[right, etiq] {$y$};
  \draw[eje] (0,0) -- (-1.3,-1.25) node[below left=-2pt, etiq] {$x$};
  \coordinate (O) at (0,0);
  \coordinate (Zp) at (0,1);
  \coordinate (Q) at (-0.75,1.05);
  \coordinate (P) at (4.6,1.8);
  \draw[vect] (O) -- ($(Q)+(0.06,-0.1)$) node[pos=0.55, left=5pt, etiq, text=white] {$\rp$};
  \cubito{Q}
  \node[etiq, font=\footnotesize, anchor=south east, text=magprof] at ($(Q)+(-0.12,0.12)$) {$\dif q$};
  \draw[vect] (O) -- (P) node[pos=0.62, below right=-2pt, etiq] {$\vc{r}$};
  \draw[vectrojo] ($(Q)+(0.11,0.02)$) -- (P) node[pos=0.62, above=2pt, etiq, text=rojoteo] {$\vc{R}=\vc{r}-\rp$};
  \pic[draw=white, line width=0.5pt, angle radius=0.45cm, "$\theta'$", angle eccentricity=1.75, font=\scriptsize, text=white] {angle=Zp--O--Q};
  \fill[magprof] (P) circle (1.5pt) node[right=2pt, etiq] {$P$};
  \draw[<-, fino] (1.2,-1.2) -- (2.0,-1.75) node[right, etiq, font=\footnotesize] {$\rho=\rho(r')$, \ $r'\le R$};
\end{tikzpicture}
\caption{Problema~2.5. Distribución esféricamente simétrica, $\rho=\rho(r')$, confinada en una esfera de radio $R$ centrada en el origen (el sombreado representa una densidad
decreciente con $r'$). Al usar coordenadas esféricas para el punto fuente, todas las integrales de los momentos se separan en un factor radial y dos angulares, y son estos
últimos los que se anulan.}
\label{fig:sem-esfera-rho}
\end{figure}


\begin{observacion}[Interpretación: el teorema de Gauss, de nuevo]\label{obs:sem-esfera-rho}
El factor polar del término cuadrupolar es $\int_{-1}^{1}\Pl{2}(u)\,\dif u=0$, y lo mismo ocurre para todo $\ell\ge1$: $\int_{-1}^{1}\Pl{\ell}(u)\,\dif u=\int_{-1}^{1}\Pl{\ell}\Pl{0}\,\dif u=0$ por la
ortogonalidad de los polinomios de Legendre. Así, \emph{todos} los momentos multipolares con $\ell\ge1$ de una distribución esféricamente simétrica se anulan, y fuera de ella el
potencial es exactamente $\phi=Q/(4\pi\eo r)$: el de una carga puntual en el centro. Es el mismo resultado que obtuvimos con la ley de Gauss en el Tema~1~\cite{TEMA1}, ahora
deducido término a término.
\end{observacion}

% =========================================================================
\subsection{Resumen de resultados}\label{sem:resumen}
% =========================================================================
\begin{table}[ht]
\centering
\caption{Resultados de los seminarios del Tema~2.}
\label{tab:resultados}
\renewcommand{\arraystretch}{1.75}
\small
\begin{tabular}{@{}>{\raggedright\arraybackslash}p{0.2\textwidth}>{\raggedright\arraybackslash}p{0.76\textwidth}@{}}
  \toprule
  \textsc{Problema} & \textsc{Resultado}\\
  \midrule
  2.1 Barra ($r>10\ell$) & $\phi\simeq\dfrac{\lambda\ell}{4\pi\eo r}+\dfrac{\lambda\ell^{3}(3\cos^{2}\theta-1)}{96\pi\eo r^{3}}$\\
                         & $\vE\simeq\Bigl[\dfrac{\lambda\ell}{4\pi\eo r^{2}}+\dfrac{\lambda\ell^{3}(3\cos^{2}\theta-1)}{32\pi\eo r^{4}}\Bigr]\ur+\dfrac{\lambda\ell^{3}\sen\theta\cos\theta}{16\pi\eo r^{4}}\,\uth$\\
  2.2 Carga en $(a,b,c)$ & $Q=q$, \ $\vp=q(a\ux+b\uy+c\uz)$, \ $Q_{xx}=q(2a^{2}-b^{2}-c^{2})$ (y cíclicos), \ $Q_{xy}=3qab$, $Q_{xz}=3qac$, $Q_{yz}=3qbc$; con $-q$ en el origen solo cambia $Q\to0$\\
  2.3 Paralelepípedo     & $\phi\simeq\dfrac{Q}{4\pi\eo r}+\dfrac{Q}{24\pi\eo r^{5}}\bigl[x^{2}(2a^{2}-b^{2}-c^{2})+y^{2}(2b^{2}-a^{2}-c^{2})+z^{2}(2c^{2}-a^{2}-b^{2})\bigr]$\\
  2.4 Esfera $\sigma_0\cos\theta$ & $\vp=\frac43\pi a^{3}\sigma_0\,\uz$ (independiente del origen, porque $Q=0$)\\
  2.5 Simetría esférica  & $\vp=\vc{0}$, \ $\Qjk=0$; en general, todos los momentos con $\ell\ge1$ son nulos\\
  \bottomrule
\end{tabular}
\end{table}
